\documentclass[11pt]{article}

\usepackage[margin=1in]{geometry}
\usepackage[T1]{fontenc}
\usepackage[utf8]{inputenc}
\usepackage{lmodern}
\usepackage{microtype}
\usepackage{amsmath,amssymb,amsthm,mathtools}
\usepackage{booktabs}
\usepackage{enumitem}
\usepackage{tabularx}
\usepackage[hidelinks]{hyperref}

\newtheorem{theorem}{Theorem}[section]
\newtheorem{proposition}[theorem]{Proposition}
\newtheorem{lemma}[theorem]{Lemma}
\newtheorem{corollary}[theorem]{Corollary}
\theoremstyle{definition}
\newtheorem{definition}[theorem]{Definition}
\newtheorem{assumption}[theorem]{Assumption}
\newtheorem{example}[theorem]{Example}
\theoremstyle{remark}
\newtheorem{remark}[theorem]{Remark}

\newcolumntype{L}[1]{>{\raggedright\arraybackslash}p{#1}}
\newcolumntype{Y}{>{\raggedright\arraybackslash}X}

\newcommand{\cA}{\mathcal A}
\newcommand{\cB}{\mathcal B}
\newcommand{\cC}{\mathcal C}
\newcommand{\cD}{\mathcal D}
\newcommand{\cF}{\mathcal F}
\newcommand{\cH}{\mathcal H}
\newcommand{\cI}{\mathcal I}
\newcommand{\cO}{\mathcal O}
\newcommand{\cP}{\mathcal P}
\newcommand{\cR}{\mathcal R}
\newcommand{\cS}{\mathcal S}
\newcommand{\cX}{\mathcal X}
\newcommand{\cY}{\mathcal Y}
\newcommand{\Id}{\operatorname{Id}}
\newcommand{\dist}{\operatorname{dist}}
\newcommand{\MTT}{Modal Triplet Theory}

\title{\vspace{-1em}
Coherence Capacity as a Control Budget\\
for Effective Descriptions:\\[0.25em]
\large Robustness, Composition, and the Limits of Resource Language}
\author{Peter Nero}
\date{Version 4, July 2026}

\begin{document}
\maketitle

\begin{abstract}
Coherence capacity is a normalized summary of how far a declared effective
construction lies from failure of one of its control conditions.  This paper
asks when that diagnostic can also be called a \emph{resource}.  The answer
requires additional operational data.  One must specify an allowed
perturbation family, a cost, an action of perturbations on states, a
composition rule, and, for a genuine resource theory, a class of free
operations and a conversion preorder.

Given those data, we define the control budget as the least perturbation
cost that exits the admissible set.  In a metric perturbation model this is
exactly the distance to inadmissibility and hence the largest guaranteed
open perturbation ball.  We prove a cumulative-consumption theorem for
normalized reserve vectors, a parallel-composition bottleneck law, and
controlled transport bounds between models.  These results make capacity
operationally useful as a robustness budget.  They also show why it is
generally neither additive nor conserved: repair, feedback, redundancy, and
changes of metric can alter the budget.

A capacity value alone is not a probability measure, force, entropy,
thermodynamic monotone, gravitational coupling, reset law, or arrow of time.
Noninjectivity also does not rule out a right inverse; an orthogonal
projection has the inclusion of its range as a bounded section.  Any
physical use therefore needs a separate action, channel, constitutive law,
instrument, state count, or causal geometry.  Within MTT, the normalized
margin paper supplies the capacity record and the Foundation supplies the
complete admissibility ledger.  The present paper contributes the
model-specific control-budget and resource-completion layer.  It does not
promote coherence capacity to a universal substance underlying all
effective physics.
\end{abstract}

\section*{Revision note: Version 4}
\addcontentsline{toc}{section}{Revision note: Version 4}
\begin{description}[leftmargin=2.5cm,style=nextline]
\item[Supersedes]
Version 4 supersedes Version 3.0, DOI
\href{https://doi.org/10.5281/zenodo.18322036}
{10.5281/zenodo.18322036}.

\item[Reason]
Version 3 called an arbitrary positive margin the fundamental resource of
effective physics.  It then used the same scalar to infer a missing right
inverse, irreversible shadow dynamics, gravity, time, entropy, measurement
collapse, horizons, thermalization, and undecidability.  None of those
conclusions follows from a zero set or stability margin alone.

\item[Resolution]
The title and central theorem have been replaced.  Capacity becomes an
operational control budget only after a perturbation class and cost are
specified.  A full resource interpretation additionally requires free
operations, composition, and conversion rules.  The paper proves
robust-radius, cumulative-consumption, composition, and transport results,
while assigning probability, force, entropy, gravity, and continuation laws
to separate source theorems.

\item[Retained content]
The physically useful intuition survives: a controlled effective
description has a finite tolerance to disturbance, and its weakest reserve
can organize robustness calculations.

\item[Remaining boundary]
MTT does not yet select one universal perturbation metric, cost, free
operation class, or conservation law for every encoding.  This paper
therefore establishes a reusable conditional framework, not a universal
physical resource theory.
\end{description}

\tableofcontents

\section{Introduction and roadmap}
\label{sec:intro}

The word \emph{resource} has several meanings in physics and mathematics.
Energy is a conserved or exchanged dynamical quantity.  Entanglement is a
resource only relative to a specified class of free operations.  A robust
control margin is a budget because it measures the smallest allowed
perturbation that destroys stability.  These meanings cannot be exchanged
by analogy.

The corrected capacity paper defines a sourced, dimensionless reserve
vector, a bottleneck scalar, and a metric distance to inadmissibility
\cite{CapacityMargin2026}.  The MTT Foundation owns the complete ledger of
conditions being summarized \cite{MTTFoundation2026}.  Those constructions
answer:
\begin{quote}
How close is this declared effective description to failure?
\end{quote}
This paper addresses the next question:
\begin{quote}
Under what additional assumptions does that clearance become a usable
control budget or a resource monotone?
\end{quote}

The answer is deliberately layered.
\begin{enumerate}[leftmargin=2em]
\item Section~\ref{sec:levels} distinguishes a diagnostic margin, an
      operational budget, a resource theory, and a physical field or charge.
\item Sections~\ref{sec:budget}--\ref{sec:composition} define perturbation
      budgets and prove their robustness, consumption, and composition
      properties.
\item Section~\ref{sec:resource} states what is still required for a genuine
      resource theory.
\item Section~\ref{sec:examples} gives concrete operator and algorithmic
      examples.
\item Section~\ref{sec:limits} identifies the former deductions that do not
      follow.
\item Sections~\ref{sec:mtt}--\ref{sec:contract} locate the framework within
      MTT and give a falsifiable completion contract.
\end{enumerate}

\section{Four levels that must not be conflated}
\label{sec:levels}

\begin{definition}[Diagnostic capacity]
A diagnostic capacity record consists of normalized signed slacks
\[
 r(x)=\bigl(r_i(x)\bigr)_{i\in\cI}
\]
for a finite list of sourced admissibility conditions.  Its signed
bottleneck and clipped scalar are
\[
 D(x)=\min_i r_i(x),
 \qquad
 C(x)=\min\{1,\max\{0,D(x)\}\}.
\]
\end{definition}

This level says which certificate is closest to failure.  It is defined by
the normalized-margin paper \cite{CapacityMargin2026}.  A value of
\(C=0.2\) has no operational meaning until the normalization and
perturbations are stated.

\begin{definition}[Operational control budget]
An operational control budget specifies what disturbances are allowed, how
they act, and how much they cost.  Its value at \(x\) is the least cost of a
disturbance that takes \(x\) outside the declared admissible set.
\end{definition}

This is analogous to a stability radius in robust control, where distance
to instability is defined relative to a structured perturbation class
\cite{HinrichsenPritchard1986}.  Different uncertainty structures can give
different radii for the same unperturbed system.

\begin{definition}[Resource theory]
A resource theory specifies objects, allowed compositions, a distinguished
class of free objects or free operations, and a conversion preorder.  A
resource monotone is a function whose allowed direction under free
operations is fixed by that theory.
\end{definition}

This usage follows the general mathematical theory of resources
\cite{CoeckeFritzSpekkens2016}.  Quantum resource theories likewise begin
with operational restrictions rather than with a scalar alone
\cite{ChitambarGour2019}.

\begin{definition}[Dynamical or physical resource]
A dynamical resource is a quantity entering a supplied action, Hamiltonian,
continuity equation, balance law, channel capacity, or constitutive
relation.  Its units, exchange law, sources, and observables are part of the
model.
\end{definition}

The implications run only downward after extra data are supplied.  A
diagnostic margin can become an operational budget.  A budget can be placed
inside a resource theory.  A resource monotone can enter a physical model.
None of those promotions is automatic.

\begin{center}
\begin{tabularx}{\textwidth}{@{}L{0.23\textwidth}Y Y@{}}
\toprule
Level & Required data & What is not yet supplied \\
\midrule
Diagnostic margin
& Slacks, units, scales, norms, tolerances, provenance
& Perturbation meaning and cost\\

Operational budget
& Allowed perturbations, action, cost, failure set
& Free operations and conversion order\\

Resource theory
& Objects, free operations, composition, monotones
& A physical conservation or exchange law\\

Physical resource
& Action or constitutive law, units, sources, observables
& Nothing further only if the complete physical model is supplied\\
\bottomrule
\end{tabularx}
\end{center}

\section{Perturbation models and the robust budget}
\label{sec:budget}

\begin{definition}[Perturbation model]
\label{def:perturbation}
A perturbation model is a tuple
\[
 \cR_{\rm ctrl}
 =
 \bigl(
 \cX,\cA,\cP,\odot,c,\cD,\pi
 \bigr),
\]
where:
\begin{enumerate}[leftmargin=2.2em]
\item \(\cX\) is the state-and-control space and
      \(\cA\subseteq\cX\) the declared admissible set;
\item \(\cP_x\) is the set of perturbations allowed at \(x\);
\item \(p\odot x\) is the perturbed state when defined;
\item \(c_x(p)\in[0,\infty]\) is its cost;
\item \(\cD\) records domains, regularity, and composition conventions; and
\item \(\pi\) records the source and units of all these choices.
\end{enumerate}
\end{definition}

\begin{definition}[Exit budget]
\label{def:exitbudget}
For \(x\in\cA\), define
\[
 B_{\cR}(x)
 =
 \inf\bigl\{
 c_x(p):p\in\cP_x,\ p\odot x\notin\cA
 \bigr\}.
\]
If no allowed perturbation exits \(\cA\), set \(B_{\cR}(x)=+\infty\).
\end{definition}

The budget is relative to \(\cP\) and \(c\).  Restricting perturbations can
increase it; changing the cost units rescales it.  These are not defects.
They are why the perturbation model belongs in every statement.

\subsection{Metric perturbations}

Suppose \((\cX,d)\) is a metric space, every \(y\in\cX\) can be treated as
the endpoint of an allowed perturbation from \(x\), and its cost is
\(c_x(y)=d(x,y)\).  Then Definition~\ref{def:exitbudget} becomes
\[
 B_d(x)=\dist_d(x,\cX\setminus\cA).
\]

\begin{theorem}[Maximal guaranteed perturbation ball]
\label{thm:ball}
Let \(\cA\) be open in \((\cX,d)\), and let \(x\in\cA\).  Then:
\begin{enumerate}[leftmargin=2em]
\item every \(y\) with \(d(x,y)<B_d(x)\) lies in \(\cA\);
\item for every \(\varepsilon>0\), there exists
      \(z\in\cX\setminus\cA\) with
      \(d(x,z)<B_d(x)+\varepsilon\), unless the complement is empty.
\end{enumerate}
Thus \(B_d(x)\) is exactly the supremal radius of an open metric ball around
\(x\) guaranteed to remain admissible.
\end{theorem}

\begin{proof}
If \(d(x,y)<B_d(x)\) and \(y\notin\cA\), then \(y\) belongs to the set over
which the defining infimum is taken, contradicting
\(B_d(x)\le d(x,y)\).  The second statement is the defining approximation
property of an infimum.
\end{proof}

This theorem supplies a precise sense in which capacity is a resource:
it is a tolerance budget for a declared class of perturbations.  It does
not say that the budget is carried by a new physical field.

\subsection{Controlled transport between equivalent models}

A coordinate or representation change need not preserve a numerical budget
exactly.  It does preserve the budget up to explicit constants when the
change controls distances in both directions.

\begin{theorem}[Bi-Lipschitz transport of exit budgets]
\label{thm:transport}
Let \((\cX,d_X)\) and \((\cY,d_Y)\) be metric spaces with admissible sets
\(\cA_X\subseteq\cX\) and \(\cA_Y\subseteq\cY\).  Suppose
\(F:\cX\to\cY\) is a bijection satisfying
\[
 F(\cA_X)=\cA_Y
\]
and, for constants \(0<m\le M<\infty\),
\[
 m\,d_X(x,z)
 \le
 d_Y(Fx,Fz)
 \le
 M\,d_X(x,z)
 \qquad(x,z\in\cX).
\]
Then the metric exit budgets obey
\[
 m\,B_X(x)
 \le
 B_Y(Fx)
 \le
 M\,B_X(x)
 \qquad(x\in\cA_X).
\]
\end{theorem}

\begin{proof}
Bijectivity and \(F(\cA_X)=\cA_Y\) imply
\(F(\cX\setminus\cA_X)=\cY\setminus\cA_Y\).  Therefore
\[
 B_Y(Fx)
 =
 \inf_{z\notin\cA_X}d_Y(Fx,Fz).
\]
Applying the two metric-comparison inequalities inside the infimum gives
the result.
\end{proof}

Exact invariance is the special case \(m=M=1\).  A one-sided Lipschitz map,
a nonbijective projection, or a map that does not carry the complete
failure set to the complete failure set does not satisfy this theorem.
Those cases require their own decoder or quotient analysis.

\subsection{Connecting the budget to normalized margins}

Let \(D(x)=\min_i r_i(x)>0\) be a normalized bottleneck.  Suppose the margin
map is \(L\)-Lipschitz in the sup norm:
\[
 \|r(y)-r(x)\|_\infty\le Ld(x,y).
\]
Then the perturbation proposition in the capacity paper gives
\[
 B_d(x)\ge\frac{D(x)}{L}.
\]
The margin value is therefore a certified \emph{lower bound} on the true
metric exit budget.  Equality needs an active perturbation reaching the
boundary at that cost.  It cannot be inferred from Lipschitz continuity
alone.

\section{Cumulative consumption and repair}
\label{sec:consumption}

An operational budget becomes especially useful for sequences of
controlled transformations.

\begin{assumption}[Rowwise consumption bound]
\label{ass:consumption}
Let \(x_0,x_1,\ldots,x_N\in\cX\), and let
\(r_i(x)\) be normalized signed reserves.  For every step \(k\), suppose a
nonnegative consumption vector
\[
 a^{(k)}=\bigl(a_i^{(k)}\bigr)_{i\in\cI}
\]
is certified such that
\[
 r_i(x_{k+1})
 \ge
 r_i(x_k)-a_i^{(k)}
 \qquad
 \text{for all }i.
\]
\end{assumption}

\begin{theorem}[Cumulative budget certificate]
\label{thm:cumulative}
Under Assumption~\ref{ass:consumption},
\[
 r_i(x_N)
 \ge
 r_i(x_0)-\sum_{k=0}^{N-1}a_i^{(k)}
 \qquad
 \text{for every }i.
\]
Hence the complete sequence remains admissible whenever
\[
 \sum_{k=0}^{N-1}a_i^{(k)}<r_i(x_0)
 \qquad
 \text{for every }i.
\]
A simpler sufficient condition is
\[
 \sum_{k=0}^{N-1}\|a^{(k)}\|_\infty<D(x_0).
\]
\end{theorem}

\begin{proof}
Iterating the rowwise inequality gives the first display.  If each final
lower bound is positive, every final reserve is positive.  For the scalar
condition,
\[
 \sum_k a_i^{(k)}
 \le
 \sum_k\|a^{(k)}\|_\infty
 <
 D(x_0)
 \le r_i(x_0)
\]
for every row \(i\).
\end{proof}

The vector condition is sharper than the scalar one because different
operations may consume different rows.  An operation that spends spectral
gap need not spend the same amount of truncation reserve.

\subsection{Repair and feedback}

The inequality in Assumption~\ref{ass:consumption} is one-sided.  A repair
or feedback operation can increase some \(r_i\).  Redundancy can move the
effective state farther from a failure set.  Re-estimating an operator in a
better norm can improve a certified lower bound without changing the
physical state.

Therefore capacity is not generally conserved.  A balance equation would
need to distinguish:
\[
 \text{consumption},\qquad
 \text{repair},\qquad
 \text{exchange},\qquad
 \text{reclassification}.
\]
The capacity definition supplies none of those rates.

\section{Composition laws}
\label{sec:composition}

Resource language is incomplete without composition.  Even in the simplest
independent case, the natural law is a bottleneck rather than a sum.

\begin{theorem}[Independent parallel bottleneck]
\label{thm:parallel}
Let \((\cX_A,d_A)\) and \((\cX_B,d_B)\) have open admissible sets
\(\cA_A\) and \(\cA_B\).  Equip
\(\cX_A\times\cX_B\) with the max-product metric
\[
 d_\infty\bigl((x_A,x_B),(y_A,y_B)\bigr)
 =
 \max\{d_A(x_A,y_A),d_B(x_B,y_B)\},
\]
and let the product admissible set be
\(\cA_A\times\cA_B\).  Then
\[
 B_\infty(x_A,x_B)
 =
 \min\{B_A(x_A),B_B(x_B)\}.
\]
\end{theorem}

\begin{proof}
The complement of the product admissible set is
\[
 \bigl((\cX_A\setminus\cA_A)\times\cX_B\bigr)
 \cup
 \bigl(\cX_A\times(\cX_B\setminus\cA_B)\bigr).
\]
The distance to the first set is \(B_A(x_A)\), because the second coordinate
may remain fixed; the distance to the second is \(B_B(x_B)\).  Distance to
their union is the minimum of the two distances.
\end{proof}

Parallel capacity is therefore limited by the weaker subsystem.  This is
not an extensive thermodynamic law.

\subsection{Coupled and redundant systems}

If the admissible set is not a Cartesian product, Theorem~\ref{thm:parallel}
does not apply.  Coupling can introduce new joint failure surfaces.
Conversely, an error-correcting code can make the logical admissible set
larger than the naive product.  A legitimate composition theorem must
specify:
\begin{enumerate}[leftmargin=2em]
\item the joint state space and metric;
\item the joint admissible set;
\item correlations in the perturbation family;
\item encoding and decoding maps;
\item resources consumed by repair; and
\item whether the reported budget is component, logical, or end-to-end.
\end{enumerate}

There is no universal formula \(B_{AB}=B_A+B_B\).  Additivity, subadditivity,
or superadditivity is a model result.

\section{When does this become a resource theory?}
\label{sec:resource}

Let \(\cS\) be a class of states and \(\cF\) a class of operations closed
under identity and composition.  Define a convertibility preorder by
\[
 x\succeq_{\cF} y
 \quad\Longleftrightarrow\quad
 y=F(x)\text{ for some }F\in\cF.
\]
A conventional resource monotone \(M\) satisfies
\[
 M(F(x))\le M(x)
 \qquad(F\in\cF).
\]
Whether larger or smaller numerical values mean ``more resource'' is a
convention, but the allowed direction must be consistent.

\begin{proposition}[Capacity is not automatically a monotone]
\label{prop:notmonotone}
An exit budget \(B\) need not be monotone under an arbitrary admissibility-
preserving operation.  In particular, a repair operation can map \(x\) to a
state \(F(x)\) with \(B(F(x))>B(x)\), while a degrading operation can produce
\(B(F(x))<B(x)\).
\end{proposition}

\begin{proof}
Take \(\cX=\mathbb R\), \(\cA=(-1,1)\), and
\(B(x)=1-|x|\).  On the admissible subdomain \((0,1)\), the map
\(F_{\rm repair}(x)=x/2\) satisfies
\[
 B(F_{\rm repair}(x))=1-\frac{x}{2}>1-x=B(x),
\]
whereas \(F_{\rm degrade}(x)=(x+1)/2\) satisfies
\[
 B(F_{\rm degrade}(x))
 =
 \frac{1-x}{2}
 <
 1-x
 =
 B(x).
\]
Both maps take \((0,1)\) into \((0,1)\).
\end{proof}

The example is elementary, but the lesson is general.  If repair is free,
capacity can be created for free and is not a resource monotone in the usual
direction.  If repair consumes ancillas, work, communication, or control
authority, those inputs must be included in the resource ledger.

\subsection{A valid model-specific monotone}

One can still define a resource theory by choosing free operations that
obey a declared comparison.  For example, suppose every \(F\in\cF\) has a
certified constant \(\alpha_F>0\) such that
\[
 B(F(x))\ge\alpha_F B(x).
\]
This is a robustness-preservation guarantee, not a monotonicity theorem.
Alternatively, choose operations satisfying \(B(F(x))\le B(x)\) and use
\(B\) as a consumable monotone.  The choice depends on the operational task.

\begin{definition}[Capacity-resource completion]
A capacity record is promoted to a capacity resource theory only when the
following are supplied:
\[
 \bigl(
 \cS,\cF,\otimes,\succeq_{\cF},
 B,\text{conversion rates},\text{catalysts},
 \text{error convention}
 \bigr).
\]
\end{definition}

This is the same discipline used in established resource theories: the
restriction on operations gives the resource its meaning
\cite{CoeckeFritzSpekkens2016,ChitambarGour2019}.

\section{Concrete conditional examples}
\label{sec:examples}

\subsection{Spectral-gap robustness}

Let \(A\) be a Hermitian matrix with a selected spectral cluster separated
from the remainder by a gap \(\gamma>0\).  Let \(E\) be a Hermitian
perturbation.  Standard spectral perturbation theory bounds eigenvalue
motion by \(\|E\|\) and invariant-subspace rotation in terms of the gap
\cite{Kato1966,DavisKahan1970}.

If the two sides of the selected gap can each move by at most \(\|E\|\),
then
\[
 \operatorname{gap}(A+E)
 \ge
 \gamma-2\|E\|.
\]
Thus
\[
 \|E\|<\frac{\gamma}{2}
\]
is a sufficient robustness budget for preserving positive separation.
It may not be the exact structured stability radius.  Restrictions on the
allowed perturbation matrix can increase the true budget.

This example makes every ingredient visible:
\begin{enumerate}[leftmargin=2em]
\item the state is the operator \(A\);
\item the perturbations are Hermitian matrices \(E\);
\item the cost is the chosen operator norm;
\item failure is closure of the selected gap; and
\item \(\gamma/2\) is a certified lower bound, not a new energy.
\end{enumerate}

\subsection{A contraction budget}

Suppose a family of maps has contraction factors \(q(x)<1\) in a declared
norm and invariant set.  The normalized reserve \(1-q(x)\) can bound how
much a perturbation may increase the Lipschitz constant before contraction
is lost.  But the invariant-set reserve, domain, and fixed-point theorem
remain separate rows.  Fixed Points I owns the corresponding conditional
MTT framework \cite{FixedPointsI2026}.

\subsection{A guarded numerical algorithm}

Let an algorithm evolve while \(C(x)>0\).  Introducing
\[
 U_\varepsilon(x)=-\log\bigl(\varepsilon+C(x)\bigr)
\]
and following \(-\nabla U_\varepsilon\) makes capacity a penalty potential
by explicit design.  The resulting force depends on the representative
\(C\), its metric, and \(\varepsilon\).  Replacing \(C\) by \(C^3\) changes
the trajectory even though the admissible set is unchanged.

At \(C=0\), the algorithm still needs a guard and one of:
\[
 \text{stop},\qquad
 \text{continue in another chart},\qquad
 \text{apply a reset map or kernel}.
\]
The capacity boundary does not choose the reset or its probabilities.

\subsection{A channel or experimental tolerance}

An experiment may define the budget as the largest calibration error,
noise strength, or control drift for which an output guarantee remains
inside tolerance.  This can be directly useful.  It is still not the
probability of an outcome.  Probabilities require a state and an
instrument; the budget only certifies the domain on which their predictions
are controlled.

\section{What the budget does not derive}
\label{sec:limits}

\subsection{No right-inverse theorem}

Let \(P:\cH\to P\cH\) be an orthogonal projection.  If
\(\ker P\ne\{0\}\), then \(P\) is noninjective.  Nevertheless the inclusion
\[
 \iota:P\cH\hookrightarrow\cH
\]
satisfies
\[
 P\circ\iota=\Id_{P\cH}.
\]
It is a bounded measurable right inverse.  Noninjectivity prevents a left
inverse that recovers every original upper state; it does not prevent a
section selecting one representative.

Capacity exhaustion may invalidate a particular decoder or commuting
diagram, but that requires a decoder theorem.  It does not follow from
noninjectivity.

\subsection{No automatic irreversibility or reset}

A first exit says that one lower certificate has ended.  The upper evolution
may remain invertible, another lower chart may exist, or the description may
simply stop.  A reset law defines a hybrid system and must separately prove
existence, measurability, conservation, and outcome probabilities.

\subsection{No probability}

A positive scalar and a boundary do not define a normalized measure.
Quadratic quantum probabilities require a state, observable algebra, and
instrument or a same-source probability theorem.  Capacity can restrict the
domain of such a theorem but cannot replace it.

\subsection{No force or equation of motion}

A gradient becomes a force only after it appears in an action, Hamiltonian,
or constitutive law with declared units and sign.  Arbitrary
reparameterization of the same zero set changes the gradient, so a
same-zero-set capacity class cannot select a unique force.

\subsection{No entropy or thermodynamics}

Entropy requires a state count, measure, density operator, coarse-graining,
or thermodynamic potential.  A control budget can constrain accessible
states without being their entropy.  A conservation or balance law requires
a symmetry or constitutive postulate.

\subsection{No gravity or Newton constant}

Postulating
\[
 S[g,C]=\int\sqrt{-g}\,f(C)R
\]
adds a scalar-tensor coupling unless \(f(C)\) is fixed and constant.
The capacity record does not select \(f\), the Einstein-Hilbert
normalization, or the observed value of \(G\).  Identifying
\(G^{-1}=C\) is also incompatible with arbitrary allowed changes of
capacity representative.

\subsection{No horizon, collapse, or arrow of time}

Horizons require Lorentzian causal geometry and field equations.
Measurement requires a system-apparatus interaction and outcome records.
An arrow of time requires asymmetric dynamics, boundary conditions,
record stability, or thermodynamic structure.  A list of first-exit events
does not supply those ingredients.

\section{Role within Modal Triplet Theory}
\label{sec:mtt}

The capacity papers now have distinct roles.

\begin{center}
\begin{tabularx}{\textwidth}{@{}L{0.28\textwidth}Y@{}}
\toprule
Source & Canonical responsibility \\
\midrule
MTT Foundation
& Complete admissibility ledger, independent control rows, and the
  stop/continue/reset alternatives \cite{MTTFoundation2026}.\\

Coherence Capacity: normalized margins
& Signed slacks, normalization scales, reserve vector, bottleneck scalar,
  metric clearance, active rows, and capacity-equivalence class
  \cite{CapacityMargin2026}.\\

This paper
& Perturbation model, exit budget, maximal guaranteed ball, cumulative
  consumption, independent composition, and resource-completion contract.\\

Fixed Points I
& Conditional spectral-projector, invariant-set, existence, contraction,
  and stability gates \cite{FixedPointsI2026}.\\

Future dynamics papers
& Any transport, repair, exchange, conservation, force, entropy, or
  geometric response law.\\
\bottomrule
\end{tabularx}
\end{center}

This hierarchy prevents a loop.  One first proves the control rows, then
normalizes them, then defines an operational perturbation budget, and only
then asks whether a task-specific resource theory or physical law exists.

\section{Completion contract and falsifiability}
\label{sec:contract}

A claim that coherence capacity is a resource for a particular model must
publish at least the following rows.

\begin{center}
\begin{tabularx}{\textwidth}{@{}L{0.25\textwidth}Y@{}}
\toprule
Row & Required content \\
\midrule
State and domain
& State space, operator domains, boundary conditions, and admissible set.\\

Capacity source
& Full normalized reserve vector, scales, tolerances, provenance, and active
  row; never the scalar alone.\\

Perturbations
& Allowed structured disturbances and their action on states.\\

Cost
& Metric or cost functional, units, uncertainty, and normalization.\\

Budget theorem
& Exact radius or certified lower and upper bounds with error control.\\

Composition
& Parallel, sequential, correlated, and redundant composition rules.\\

Free operations
& Operations regarded as costless and proof that they preserve the required
  closure properties.\\

Monotone direction
& Whether capacity must increase, decrease, or remain comparable under free
  operations.\\

Repair ledger
& Ancillas, work, communication, feedback, or other resources consumed by
  stabilization.\\

Physical promotion
& Separate action, instrument, measure, entropy, or geometry theorem if a
  physical interpretation is claimed.\\

Provenance
& Source-independent inputs, fitted inputs, held-out outputs, intervals,
  code, commit, and hashes for numerical values.\\
\bottomrule
\end{tabularx}
\end{center}

The framework is falsifiable at several levels.  A claimed radius fails if
an allowed cheaper perturbation exits the admissible set.  A claimed
composition law fails if a joint perturbation violates its bound.  A
resource monotone fails if a declared free operation moves it in the
forbidden direction.  A physical promotion fails if the predicted
observable or held-out comparison fails.  These are stronger tests than
asking whether a scalar happened to vanish at a chosen boundary.

\section{Conclusion}
\label{sec:conclusion}

Coherence capacity can be a genuine \emph{control budget}: after the
perturbations and their cost are fixed, it measures the least disturbance
that destroys a declared effective description.  That interpretation gives
precise robustness, cumulative-consumption, and composition theorems.

It is not, by definition, the fundamental resource of all effective
physics.  A resource theory needs operational restrictions and conversion
rules.  A physical resource needs dynamics, units, sources, and observables.
Keeping those levels separate makes the capacity idea more useful, because
each proposed application now has an exact completion contract and an
observable failure mode.

\section*{Reproducibility statement}

This paper reports structural definitions and proofs, not a fitted
parameter or numerical physical prediction.  The source TeX, bibliography,
revision audit, and generated PDF are versioned in the MTT papers
repository.  A future numerical control-budget result must additionally
archive its perturbation family, cost, full reserve vector, active rows,
interval bounds, code commit, and verification hashes.

\bibliographystyle{plain}
\bibliography{main}

\end{document}
